\documentclass[10pt,a4paper]{amsart}
\usepackage[margin=19mm]{geometry}
\usepackage{amsmath,amssymb,mathtools}
\usepackage[scr=rsfs,cal=euler]{mathalfa}
\usepackage{xfrac}
\usepackage[unicode,hidelinks,bookmarks=false]{hyperref}
\newcommand{\ie}{i.e.\ }
\newcommand{\cf}{cf.\ }
\newcommand{\resp}{respectively\ }
\newcommand{\Subsection}{\subsection}
\providecommand{\nobreakdash}{\nobreak\hbox{-}}
\setlength{\emergencystretch}{2em}
\multlinegap0pt
\usepackage{amssymb}
\usepackage[shortlabels]{enumitem}
\usepackage{braket}
\usepackage{bm}
\newtheorem{lemma}{Lemma}
\DeclareMathOperator{\Div}{div}
\newcommand{\IV}{\mathcal{IV}}
\newcommand{\R}{\mathbb{R}}
\DeclareMathOperator{\reg}{reg}
\DeclareMathOperator{\spt}{spt}
\title{Green functions on stationary varifolds}
\author{Yifan Guo}
\date{Source: arXiv:2404.00119v2 (CC BY 4.0)}
\begin{document}
\maketitle

\noindent\emph{Source and licence.} Green functions on stationary varifolds. Version arXiv:2404.00119v2 (CC BY 4.0), distributed by arXiv under the Creative Commons Attribution 4.0 International licence, http://creativecommons.org/licenses/by/4.0/, as recorded in the arXiv licence metadata for this exact version and shown on the abstract page https://arxiv.org/abs/2404.00119v2. Full licence notice: \texttt{licenses/arxiv-2404.00119-CC-BY-4.0.txt}.\par\smallskip

\noindent\textit{Editorial scope.} Counted unit: the article's lemma super (source0.tex lines 1470--1476). The article's complete original proof of it is delivered with it (source0.tex lines 1477--1497, one \texttt{proof} environment of 21 lines). No mathematics is written, repaired or re-derived: the statement and the proof are the article's own lines, in the article's own notation, and no retained line is substituted or re-flowed.  No further source-local context is needed: every cross-reference of every retained line resolves inside the two delivered spans.  Nothing was omitted from any retained span, so the package carries no omission record.  No retained line contains a citation, so the wrapper carries no bibliography and no bibliography entry.  No retained line contains a figure environment, an image-inclusion command or drawing code, so no image is delivered and the package declares no asset dependency.  Editorial formatting only, in addition: an AMS \texttt{amsart} preamble that restates the article's own declarations and macros used by the retained lines, namely the environments \texttt{lemma} on separate counters, together with the standard packages those lines need.  The retained environments print in this document as Lemma 1 in order of appearance, whereas the article prints the same items with the numbering of its own full document.  The complete native source of the article as distributed in the arXiv e-print for this exact version is bundled unchanged as \texttt{sources/156932/source0.tex}.
\begin{lemma}\label{super}
	Let $V\in \IV_n(U)$ be stationary. Then for any $x\in \R^{n+k}$, $\phi\in C_c^\infty(U)$ with $\phi\ge0$ and $x\notin \spt \phi$, we have
	\[\int \nabla^V\Gamma(x,\cdot)\cdot \nabla^V\phi d\|V\|=\frac{1}{\omega_n}\int \phi (y)\frac{|(y-x)^\perp|^2}{|y-x|^{n+2}}d\|V\|(y)\ge0. \]
	 Moreover, if $y\in\reg V$, then we have
	 \[\Delta_{V,y} \Gamma(x,y)=-\frac{1}{\omega_n}\frac{|(y-x)^\perp|^2}{|y-x|^{n+2}}\]
	where $(y-x)^{\perp}$ is the projection of $y-x$ to $T_yV^\perp$.
\end{lemma}

\begin{proof}
	Without loss of generality, we take $x=0$ and denote $r=|y|$. We compute
	\begin{equation}\label{div.r^2-n}
		\begin{aligned}
			\Div _V(\nabla r^{2-n})=&(2-n)\Div _V(r^{-n}x)=n(2-n)r^{-n}+(2-n)(-n)r^{-n}|\nabla ^Vr|^2\\
			=&n(2-n)r^{-n}|\nabla^\perp r|^2\\
			=&n(2-n)\frac{|y^\perp|^2}{|y|^{n+2}}.
		\end{aligned}
	\end{equation}
	For any $\phi\in C_c^\infty(U)$ with $\phi\ge0$ and $0\notin \spt \phi$, the function $r^{2-n}$ is smooth in a neighborhood of $\spt\phi$. Since $V$ is stationary, we have
	\[0=\int \Div _V(\phi\nabla r^{2-n})d\|V\|=\int \nabla^V r^{2-n}\cdot\nabla^V\phi d\|V\|+\int \phi \Div_V (\nabla r^{2-n})d\|V\|.\]
	Thus we get
	\[\int \nabla^V r^{2-n}\cdot\nabla^V\phi d\|V\|=n(n-2)\int \phi \frac{|y^\perp|^2}{|y|^{n+2}}d\|V\|\ge0.\]
	If $y\in \reg V$, we can take $\{e_i\}_{i=1}^n$ to be an orthonormal frame of $TV$ near $y$. Stationarity implies the mean curvature $\sum_{i=1}^n(\nabla_{e_i}e_i)^\perp=0$ near regular points. Hence
	\[\begin{aligned}
		\Div _V(\nabla r^{2-n})=&\Div _V(\nabla^V r^{2-n})+\Div _V(\nabla^\perp r^{2-n})=\Delta_Vr^{2-n}+\sum_{i=1}^n \nabla_{e_i}\nabla^\perp r^{2-n}\cdot e_i\\
		=&\Delta_Vr^{2-n}-\sum_{i=1}^n \nabla^\perp r^{2-n}\cdot \nabla_{e_i}e_i\\
		=&\Delta_Vr^{2-n}.
	\end{aligned}\]
	Combining with (\ref{div.r^2-n}), we get the second equation.
\end{proof}

\end{document}
